\section{研究背景及问题重述}
\subsection{研究背景}
低空飞行活动通常位于受地表影响较强的大气层内。地形起伏、建筑物阻挡和地表加热共同改变局地风场，使飞行器在较短航程内经历不同的气流状态。当气流变化的空间尺度与飞行器尺度相近时，飞行稳定性受到直接影响。因此，研究低空湍流既需要描述气象要素的分布，也需要把这些要素转化为与航路相关的风险信息。

现有探测设备提供了不同侧面的观测。地面自动站反映近地面条件，风廓线雷达提供垂直风场，微波辐射计补充温度廓线，多普勒天气雷达提供一定空间范围内的径向速度与谱宽信息。由于采样时间、探测体积和可观测变量不同，直接拼接这些数据会混合不同尺度的变化。本文先建立统一的变量定义与质量控制流程，再根据各类设备的观测能力构造湍流指标。

本研究围绕低空飞行中的风险识别需求，按照单站湍流诊断、区域三维重构和预报航路规划的顺序展开。图\ref{fig:background}给出了观测与应用之间的关系，表\ref{tab:background}列出了各类数据在模型中的作用。该结构使物理参数的计算、空间信息的融合与最终决策保持对应。
\auxfigure{background}{低空观测与航路应用关系}{fig:background}
\begin{table}[H]\centering\auxtablecaption{观测资料与研究任务的对应关系}\label{tab:background}
\begin{tabularx}{\linewidth}{lYY}\toprule
资料类别 & 可用信息 & 建模作用\\\midrule
风廓线雷达 & 风速分量、径向速度、谱宽与质量标志 & 构造垂直廓线与动力诊断量\\
微波辐射计 & 温度或热力廓线及质量标志 & 计算位温与稳定度\\
地面自动站 & 近地面风、温度等题目提供要素 & 约束低层条件与水平覆盖\\
天气雷达 & 径向速度、谱宽及站点扫描信息 & 补充三维空间结构\\
数值天气预报 & 风场、热力场及有效预报时刻 & 诊断未来湍流风险\\\bottomrule
\end{tabularx}\end{table}
